\begin{lemma}[Extension-base Part I(4): colors in one support class]
Let the active support class \(L(w)\) be nonempty and write its exact
increasing enumeration as
\[
             i_1<i_2<\cdots<i_{p+1}.
\]
Assume this is also the exact enumeration, through the terminal round, of
all PE slots supported at \(w\).  Write the certified gamma and delta fields
at slot \(i_{a+1}\) as \(\gamma_a,\delta_a\).  Suppose the conclusion of
Part I(2) is available in the form
\[
                 \delta_b\ne\gamma_a\qquad(0\leq a\leq b\leq p).
\]
Then
\[
 \gamma_{a+1}=\delta_a\qquad(0\leq a<p),
\]
and the \(p+2\) colors
\[
 \gamma_0,\ \delta_0=\gamma_1,\ \delta_1=\gamma_2,
       \ldots,\ \delta_{p-1}=\gamma_p,\ \delta_p
\]
are pairwise distinct.  Equivalently, the function on
\(\operatorname{Fin}(p+2)\) obtained by adjoining \(\gamma_0\) in front
of the delta sequence is injective.
\end{lemma}
\begin{proof}
The range condition for the displayed indices is the support-class
definition \noderef{ExtensionBaseSupportClass}; the second exact-range
hypothesis says that there are no omitted PEs supported at \(w\).  Apply
\noderef{SupportingPEColorChain} to this increasing enumeration.  It gives
\[
                    \gamma_{a+1}=\delta_a
                    \quad(0\leq a<p),                         \tag{1}
\]
using the exact gamma and delta fields from the hypotheses.

Define \(c_0=\gamma_0\) and \(c_{b+1}=\delta_b\) for
\(0\leq b\leq p\).  We prove directly that the \(c_r\) are pairwise
distinct.  First, for every \(b\), Part I(2) with \(a=0\) gives
\(c_{b+1}=\delta_b\ne\gamma_0=c_0\).

Now take \(0\leq a<b\leq p\).  Since \(a+1\leq b\), Part I(2), with
\(a+1\) in its first position and \(b\) in its second, gives
\[
                         \delta_b\ne\gamma_{a+1}.              \tag{2}
\]
Equation (1) identifies \(\gamma_{a+1}=\delta_a\).  Hence (2) says
\(c_{b+1}=\delta_b\ne\delta_a=c_{a+1}\).  Together with the preceding
comparison with \(c_0\), this handles every pair of distinct indices;
reversing the pair handles the opposite order.  Thus the finite function
\(c=\operatorname{Fin.cons}(\gamma_0,\delta)\) is injective, which is
exactly the asserted distinctness and completes Part I(4).
\end{proof}
